\section{Introduction}
From biological organisms to engineered feedback loops, complex systems science describes how structures maintain coherence against entropic disturbances \cite{schrodinger_what_1992}. Their defining feature is \textit{persistence}; they retain their structure despite ongoing noise and decay. Control theory describes how systems regulate against fluctuations, taking the regulation target (setpoint) as an externally specified boundary condition. Information theory quantifies the cost of information processing, taking the processor as given. Thermodynamics establishes the arrow of time, treating entropy as a property of states. These three fields developed independently to solve distinct engineering problems: thermodynamics from heat engines, information theory from telecommunication, and control theory from servomechanisms. Each is optimizing a system that already exists.

We define this missing foundation: the minimal architecture presupposed by any dynamical or optimization discipline. Grounding this investigation in algorithmic information theory, we identify the invariant foundation beneath information theory itself; thermodynamics and control theory emerge as substrate-specific instantiations of the same universal mechanics. By uniting this static architecture with the minimal dynamic Action required to maintain it, this paper establishes \textit{persistence mechanics}: the necessary and sufficient conditions for a single persistent system, its dynamics, and the selection principle governing observable history.

Previous complex systems research has converged on pieces of this puzzle from multiple directions. Prigogine established that persistent structures occupy a minimum-entropy-production attractor\cite{prigogine_introduction_1961}; Friston formulated persistence as the minimization of variational free energy via an internal model \cite{friston_free-energy_2010}; England showed how external drive selects for structural configurations in non-equilibrium steady states\cite{england_dissipative_2015}; and cybernetics has catalogued the organizational patterns, such as requisite variety and good regulation, that correlate with survival\cite{ashby_introduction_1956,conant_every_1970}. While individual domains have named components performing these functions, none formally define the minimal conditions required for a persistent structure to exist in the first place.

We describe two universal mechanics found in systems that persist. First, the \textbf{Architecture of Persistence} (\Cref{sec:architecture_of_persistence}): six substrate-independent pillars extracted from the logical requirements of Finiteness, Conservation, and Causality. This architecture specifies what must exist, not how it behaves under transition. Second, the \textbf{Dynamics of Persistence} (\Cref{sec:dynamics_of_persistence}): the universal behavior governing how any persistent system dissipates, accumulates entropic action, and undergoes selection. We derive the dissipation density $\mathcal{D} = \xi \cdot \mathbf{K} \cdot f$ as the minimum dynamical description any state transition law must presuppose.

From these two components \Cref{sec:minimum_action} shows that the Principle of Least Action emerges as the boundary of survivorship: trajectories that do not minimize their Entropic Action undergo exponential suppression, leaving only the least-action remainder as observable history. \Cref{sec:open_systems} extends the architecture to open systems and \cref{sec:universality_mapping} then illustrates the architecture's mapping across biological, ecological, and computational substrates.